\section{Introduction}
\label{sec:b-introduction}

An autonomous experimentation agent that can call a verifier faces a timing
question before it faces a ranking question.  A verification budget spent at an
early checkpoint may repair a candidate that the frozen continuation would have
carried to failure; the same budget spent at a later checkpoint may act on
information that does not exist yet.  Spending early can therefore improve a
result, consume resources without benefit, or remove a more valuable later
opportunity \citep{calibrationnotcontrol2026,harnessr12026}.  This paper studies
the timing question directly for a fixed verification-and-response regime that
may be used \emph{at most once} across two online checkpoints.

\paragraph{Problem.}
Along each \emph{original root} (one independent run of the frozen agent from a
task template) two online checkpoints are eligible.  A controller may spend its
single verification at checkpoint 1, reserve it and spend it at checkpoint 2 if
that checkpoint is reached, or never verify.  Early termination and
infeasibility are part of the problem: a root that ends before checkpoint 2, or
at which verification is unavailable, remains in the original-root population
and is analysed under its registered fallback.  The regime is fixed and frozen:
the verifier checks consistency against \emph{visible} development diagnostics,
and the response rule replaces a candidate only on a strict visible improvement,
leaving ties unchanged.  The hidden evaluation is never visible to the agent or
the verifier.

\paragraph{Primary hypothesis (N2).}
The frozen two-checkpoint controller improves the original-root mean utility $U$
against \emph{both} a repeated-myopic controller and a BAVAR-inspired
context-dependent VOI scheduler, at the same allowed verifier and total-work
budgets.  The two
primary contrasts are tested with simultaneous intervals, with template families
as the sampling units and a prespecified treatment of batch coupling
(\cref{sec:b-statistics}).  Decision rules follow the registered design: a
contrast is supported when its adjusted lower confidence bound exceeds zero; the
practical target is $0.02\,U$; a contrast whose adjusted upper bound is below
$0.02$ rejects the practical target; an interval spanning the target is
inconclusive, with no automatic budget expansion.

\paragraph{Validity constraint (N1).}
Policy assessment separates sampling uncertainty from an explicitly conditional
allowance for outcome-relevant reconstruction error.  Under a declared joint
discrepancy event of probability at most $p$, bounded utility, a fixed policy
and the required calibration independence, the population gain lower bound loses
at most $2p$; this is an application of the measurement bound of the anonymous
companion paper \citep{companion2026}, not a new safety theorem
(\cref{sec:b-theory}).  There is no
pointwise gate from a global $p$: a finite frozen policy family instead receives
simultaneous population lower bounds and abstains if none qualifies.  For the
two-checkpoint comparison the allowance is a joint path allowance $p_{\text{path}}$
covering the decision histories and the compared regime outcomes.

\paragraph{Falsifiers.}
N2 is falsified if the adjusted lower confidence bound does not exceed zero for
either primary comparator, or if the adjusted upper bound falls below the
$0.02$ practical target; an interval that spans the target is inconclusive.
N1's mathematical claim is falsified by a counterexample satisfying its stated
assumptions.  Its practical claim is falsified if adding the allowance does not
improve the declared utility/harm tradeoff at equal allowed resources, including
the always-abstain controller that satisfies validity while failing the
practical contribution.  Understated $p$ is a sensitivity failure, not a
refutation of the conditional proposition.

\paragraph{What this paper is not.}
It is not a claim that prediction and intervention value come apart; that is
prior work \citep{interventionparadox2026,calibrationnotcontrol2026}.  It is not
a new general identification theorem for replay; the measurement preconditions
are developed by the anonymous companion and assumed here as a contract.  It
does not claim the historical ordinal-sufficiency identity as its contribution;
with identical ridge features, penalties and rows the advantage learner and the
action-value difference are the same estimator, and that comparison is retained
only as an implementation equivalence check
(\cref{sec:b-theory,sec:b-equivalence}).  It reports no live result: every
performance number is marked \emph{planned}
(\cref{sec:b-experiments}).

\paragraph{Organisation.}
\Cref{sec:b-problem} states the two-checkpoint problem and the original-root
population.  \Cref{sec:b-regimes} fixes the three regimes on one tree.
\Cref{sec:b-data} gives the branch-record schema and the runtime qualification
boundary.  \Cref{sec:policy} defines the matched controllers and total-work
accounting.  \Cref{sec:b-theory} states the scoped population proposition and
the equivalence check.  \Cref{sec:b-experiments} is the planned evaluation with
the candidate-consistency task.  \Cref{sec:b-statistics} fixes outcomes,
allocation estimands and inference.  \Cref{sec:b-limitations} states
limitations, \cref{sec:related} positions the work at method level, and
\cref{sec:b-bridge} records the deferral of unrestricted sequential control.
